%%%%%%%%%%%%%%%%%%%%%%%%%%%%%%%%%%%%%%%%%
% Medium Length Professional CV
% LaTeX Template
% Version 2.0 (8/5/13)
%
% This template has been downloaded from:
% http://www.LaTeXTemplates.com
%
% Original author:
% Trey Hunner (http://www.treyhunner.com/)
%
% Important note:
% This template requires the resume.cls file to be in the same directory as the
% .tex file. The resume.cls file provides the resume style used for structuring the
% document.
%
%%%%%%%%%%%%%%%%%%%%%%%%%%%%%%%%%%%%%%%%%

%----------------------------------------------------------------------------------------
%	PACKAGES AND OTHER DOCUMENT CONFIGURATIONS
%----------------------------------------------------------------------------------------


\documentclass{resume} % Use the custom resume.cls style

\usepackage{verbatim}
\usepackage{enumitem}
\usepackage{changepage} % for adjustwidth
\usepackage{url}
\usepackage{xcolor}

%\usepackage[colorlinks=true, linkcolor=blue, urlcolor=blue, citecolor=blue]{hyperref}
%\usepackage[colorlinks=true, urlcolor=black, linkcolor=black, citecolor=black]{hyperref}
\usepackage[colorlinks=true, urlcolor=blue!45!black, linkcolor=blue!50!black]{hyperref}
\usepackage{soul}





\usepackage[left=0.4 in,top=0.4in,right=0.4 in,bottom=0.4in]{geometry} % Document margins
\newcommand{\tab}[1]{\hspace{.2667\textwidth}\rlap{#1}} 
\newcommand{\itab}[1]{\hspace{0em}\rlap{#1}}
\name{\href{https://ryannnice.github.io/}{Renyuan Liu}} % Your name
%\address{123 Pleasant Lane \\ City, State 12345} % Your secondary addess (optional)
\address{+86 14784206312 \\ renyvannnn@gmail.com}  % Your phone number and email
%\address{https://ryannnice.github.io}  % Your phone number and email
\begin{document}


\begin{comment}
%----------------------------------------------------------------------------------------
%	OBJECTIVE
%----------------------------------------------------------------------------------------
\begin{rSection}{OBJECTIVE}

{Multidisciplinary engineer with excellent problem solving and critical thinking abilities, seeks hands on experience within a company that embraces creativity and innovation.Through my studies, I have gained extensive knowledge of production and manufacturing engineering, product design, among many other components of  Mechanical Engineering. Effective communicator who builds positive, cohesive relationship with all level of staff, eager to put my engineering knowledge to practical, applied use.}

\end{rSection}
\end{comment}


%----------------------------------------------------------------------------------------
%	Education
%----------------------------------------------------------------------------------------

\begin{rSection}{Education}
	
\begin{itemize}[left = -0.22in]
	\item {\bf \href{https://cuhk.edu.cn/en}{The Chinese University of Hong Kong, Shenzhen}} 
	\hfill{Sept. 2026 -- Jun. 2028 (Expected)}
	\\
	\textbf{MPhil in Artificial Intelligence}; Supervisor: \href{https://sai.cuhk.edu.cn/teacher/153}{Prof. Junjie Hu}
	
	\item {Guangzhou University} 
	%\hfill{Guangzhou, China}
	\hfill{Sept. 2022 -- Jun. 2026 (Expected)}
	\\
	B.Eng. in Information Security; {GPA: \textbf{89.81}/100.00; Ranking: \textbf{Top 8\%} }\\
	%B.Eng. - Information Security; {\textbf{GPA: 90.13/100.00}; \textbf{Ranking: Top 10\%} }\\
	Curriculum: Machine Learning 100, Operating System 98, Data Structure and Algorithm 97.
	 	
\end{itemize}

\end{rSection}


%----------------------------------------------------------------------------------------
%	Publications
%----------------------------------------------------------------------------------------

%\begin{rSection}{Publications}


\begin{rSection}{Papers}
\begin{itemize}[left = -0.22in]
	
	\item \textbf{R. Liu} and Q. Fu, 
	\href{https://ieeexplore.ieee.org/document/11227781}
	{Attention-Driven LPLC2 Neural Ensemble Model for Multi-Target Looming Detection and Localization}. 
	\textit{The 2025 International Joint Conference on Neural Networks (IJCNN)(CCF-C, AR $\approx$ 38\%)}.
	
	
	\item G. Gao\textsuperscript{*}, \textbf{R. Liu}, M. Wang and Q. Fu\textsuperscript{*}, \href{https://www.mdpi.com/2313-7673/9/11/650}
	{A Computationally Efficient Neuronal Model for Collision Detection with Contrast Polarity-Specific Feed-Forward Inhibition}. 
	\textit{Biomimetics, vol.9, no.11, p.650, 2024 (JCR Q1, IF = 3.9).}	
	
	
	\item C. Fang\textsuperscript{*}, H. Zhou, \textbf{R. Liu}, and Q. Fu\textsuperscript{*},
	\href{https://www.sciencedirect.com/science/article/pii/S092523122502332X}
	{A Neuromorphic Binocular Framework Fusing Directional and Depth Motion Cues Towards Precise Collision Prediction}. 
	\textit{Neurocomputing, 131660 (JCR Q1, CCF-C, IF = 6.5).}
	
		
	\item H. Zhou, C. Fang, \textbf{R. Liu}, and Q. Fu,
	\href{https://www.ejournal.org.cn/thesisDetails#10.12263/DZXB.20250337&lang=en}
	{A Bio-Plausible Neural Network Integrating Motion and Disparity Pathways for Looming Perception}. 
	\textit{Acta Electronica Sinica, p.1-16, 2025 (CCF-A in Chinese Category).}
	
	
	\item J. Huang\textsuperscript{*}, Z. Qin, M. Wang, \textbf{R. Liu}, and Q. Fu\textsuperscript{*}, 
	\href{https://ryannnice.github.io/assets/A biomimetic collision detection visual neural model coordinating self-and-lateral inhibitions.pdf}
	{A Biomimetic Collision Detection Visual Neural Model Coordinating Self-and-Lateral Inhibitions}.$\,$
	\textit{The$\,$14th International$\,$Conference$\,$on$\,$Biomimetic$\,$and$\,$Biohybrid Systems (Oral).} %(Living Machines 2025)}. 


	\item \textbf{R. Liu}, H. Zhou, C. Fang and Q. Fu, \href{https://arxiv.org/pdf/2509.13827}{How Fly Neural Perception Mechanisms Enhance Visuomotor Control of Micro Robots}. 
	%[Manuscript under double-blind review].
	\textit{Under review in 2026 International Conference on Robotics and Automation (ICRA)(CCF-B)}.
		
	\item M. Wang\textsuperscript{*}, \textbf{R. Liu\textsuperscript{*}} and Q. Fu,
	\href{https://ryannnice.github.io/assets/TCDS-2025-0806_Proof_hi.pdf}{Enhancing Collision-Selectivity in Autonomous Micro-Robots by Elevated Temporal Derivatives in Neuronal Assembly Framework}. 
	%[Manuscript under double-blind review]
	\textit{Under review in IEEE Transactions on Cognitive and Developmental Systems (JCR Q1, IF = 4.9).}

\end{itemize}
\end{rSection}










%----------------------------------------------------------------------------------------
%	Honors and Awards
%----------------------------------------------------------------------------------------

\begin{rSection}{Honors and Awards}
	
	\begin{itemize}[left = -0.22in]
		
		%\item 
		%\textbf{First-Class Scholarship (\textbf{Top 5\%})}, \textit{Guangzhou University} 
		%\hfill{Sept. 2025}
		
		\item
		\textbf{First Prize (Provincial; Top 3\%)},
		Chinese Collegiate Computing Competition (4C)
		\hfill{May 2025}
		
		\item 
		\textbf{Honorable Mention (International)}, 
		%\textbf{International Second Prize}, 
		Mathematical Contest in Modeling
		(MCM)
		\hfill{Jan. 2025}
		
		\item 
		\textbf{First Prize (\textbf{National; Top 5\%})}, Asia and Pacific Mathematical Contest in Modeling
		(APMCM) 
		\hfill{Nov. 2024}
		
		\item 
		\textbf{First Prize \& Innovation Silver Award (\textbf{Provincial; \textbf{Top 2} out of \textbf{1,167} Teams})},
		\hfill{Nov. 2024}\\
		“Greater$\,$Bay$\,$Area$\,$Cup”$\,$Guangdong-Hong$\,$Kong-Macao$\;$Financial Mathematics$\,$Modeling$\,$Competition
		
		\item
		\textbf{Commendation Letter} for Outstanding Performance in the Winning Team,
		\hfill{Nov. 2023}\\
		Interdisciplinary Programme at University of Hong Kong and The University of Macao
		
		\item 
		\textbf{Second; Third; First-Class Scholarship (\textbf{Top 8; 12; 5\%})}, \textit{Guangzhou University} 
		\hfill{Dec. 2025; 2024; 2023}
		

	\end{itemize}
	
\end{rSection} 





%----------------------------------------------------------------------------------------
%	Skills
%----------------------------------------------------------------------------------------
\begin{rSection}{Skills}
	
	\begin{itemize}[left = -0.22in]
		\item \textbf{Coding:} C/C++, Python, Matlab, TypeScript 
		\quad \textbf{English:} IELTS 6.5, CET-6 564
		\item \textbf{Skills:} Claude Code/Codex/OpenCode, Transformer, Ascend 910B, Agent skills, Docker, FastAPI, PyTorch, Hugging Face, Git, ROS, STM32, Webots, Keil, MS Office, MS Visio, Adobe PS/Pr
		%\item \textbf{Languages:} IELTS 6.5 (R8.0, L6.5, W6.0, S5.5), CET-6 564 (R242/248.5)
	\end{itemize}
	
\end{rSection} 




\newpage
\begin{rSection}{Internship Experience}
	
% Shenzhen Research Institute of Big Data
\begin{adjustwidth}{-0.22in}{0in}
	\textbf{\href{https://www.sribd.cn/en}{Shenzhen Research Institute of Big Data (SRIBD)}}, Shenzhen, China  
	\\
	Role: \textbf{Algorithm Group Intern}
	\quad Advisors: Jianghua Wu and Akang Wang 
	\hfill Mar. 2026 -- Present
\end{adjustwidth}
\begin{itemize}[left = -0.22in]
	\item \textbf{AI agent application development:} Developed a web-based AI education platform from scratch, including agent-driven news screening, generation and push modules, learning modules, and a coding agent creative workshop; also developed the supporting Android application.
	
	\vspace{-0.5em}
	\item \textbf{Large model algorithms:} Deployed Qwen models on the \textbf{Ascend 910B} accelerator and optimized inference speed; developed an LLM routing system whose semantic-based algorithm reduced cost by \textbf{62\%} while controlling the accuracy loss within \textbf{10\%}; currently continuing work on large-model deployment and optimization.
\end{itemize}

\end{rSection}



%----------------------------------------------------------------------------------------
%	Research Experience
%----------------------------------------------------------------------------------------
\begin{rSection}{Research Experience}
% The University of York
\begin{adjustwidth}{-0.22in}{0in}
	\textbf{Computational Autonomous Learning Systems Lab}
	\quad Advisor:  \href{https://scholar.google.com/citations?user=VxmFgc0AAAAJ&hl=zh-CN}
	{Prof. Pengcheng Liu} \\
	\href{https://www.york.ac.uk/computer-science/research/}
	{Department of Computer Science, University of York}, York, UK (On-Site)
	%\hfill \href{https://ryannnice.github.io/assets/Invitation Letter.pdf}{Jun. 2025 -- Aug. 2025 (Ongoing)} 
\end{adjustwidth}
\begin{itemize}[left = -0.22in]
	
	\item \textbf{Biologically Plausible Mechanisms for Long-Term Motion Learning}. 
	\hfill Jun. 2025 -- Sept. 2025
	\begin{itemize}[label=\textopenbullet] % 使用空心点
		\vspace{-0.2em}
		\item \textbf{Robotic arm motion planning:} Learned expert-guided trajectory optimization via Learning from Demonstration \textbf{(LfD)} and applied biologically inspired probabilistic movement primitives \textbf{(ProMPs)} for push–grasping with the \href{https://franka.de/}{\textit{Franka Emika Panda}}.
		
		\vspace{-0.5em}
		\item \textbf{Navigation and manipulation:} Working on developing a \textbf{lifelong learning} navigation–manipulation system on \href{https://www.turtlebot.com/turtlebot3/}{\textit{TurtleBot 3}} with \href{https://emanual.robotis.com/docs/en/platform/openmanipulator_x/overview/}{\textit{OpenMANIPULATOR-X}} that adapts to new environments while \textbf{retaining} performance in previously learned ones.
	\end{itemize}
\end{itemize}

% Guangzhou University
\begin{adjustwidth}{-0.22in}{0in}
	\textbf{Machine Life and Intelligence Research Centre}
	\quad Advisor:  \href{https://scholar.google.com/citations?user=YIte1M8AAAAJ\&hl=zh-CN}{Prof. Qinbing Fu} \\
	School of Mathematics and Information Sciences, Guangzhou University, Guangzhou, China
\end{adjustwidth}

\begin{itemize}[left = -0.22in]
	
	
	\item \textbf{Real-time Visual Processing Systems Development of Micro-Mobile Robot}\hfill Mar. 2023 -- Present
	\begin{itemize}[label=\textopenbullet] % 使用空心点

		%\item Reading and giving reports of research articles during research seminars on a weekly basis.%Regularly analyzed and reported on cutting-edge research papers in weekly lab seminars.
		
		\item 
		Developed the multi-attention LPLC2 (mLPLC2) neural network model inspired by the visual system of the fly by leveraging a \textbf{bottom-up attention} mechanism driven by motion-sensitive neural pathways
		(\href{https://github.com/Ryannnice/Offline_Multi-Attention_LPLC2_Model/blob/main/LPLC2.cpp}{\textbf{independently, 3k lines of code in C/C++}}).
		
		\vspace{-0.3em}
		
		\item 
		Deployed visual neural network models inspired by insect neurons onto the \href{https://link.springer.com/chapter/10.1007/978-3-319-96728-8_17}{STM32-based micro-robot \textit{Colias}}, achieving real-time collision perception and avoidance.
		Optimized model memory usage to fit within the \textbf{62 KByte} SRAM capacity of \href{https://link.springer.com/chapter/10.1007/978-3-319-96728-8_17}{\textit{Colias}};
		developed and refined algorithms to enable real-time execution under extreme computational constraints (processing time \textbf{$<$ 33 ms} on the STM32F427 micro chip);
		%in the extremely computing resource-constrained embedded system;
		performed debugging, tuning, and conducted both offline and online experiments.
		
		\vspace{-0.3em}
		\item 
		\href{https://ryannnice.github.io/assets/TAROS_2025_Poster_100.pdf}
		{Fly-Inspired Ultra-selective Looming Perception and Avoidance on Resource-Constrained Micro-Robots}, poster presentation at \textit{the 26th Towards Autonomous Robotic Systems (TAROS 2025) Conference}.
		
		\vspace{-0.3em}
		\item 
		Selected code can be accessed below: \\
		\href{https://github.com/Ryannnice/neuro-life-project/blob/main/micro_embodied/colias_core/coliasSense_LPLC2.c}
		%{https://github.com/Ryannnice/mLPLC2_Colias_Robot}
		{Fly Visuomotor-Inspired Attention-LPLC2 Model \textbf{(independently, 2k lines of code in C)}}; \\
		\href{https://github.com/Ryannnice/Optimized-LGMD/blob/main/coliasSense_LGMD.c}
		{Locust Vision-Inspired Optimized-LGMD Model \textbf{(independently, 1k lines of code in C)}}.
		
		
	\end{itemize}
	
	
\end{itemize}



%\vspace{1em}

\end{rSection} %Research Experience Over ... 








%Hobbies: Movie, Music, Photography, Basketball, Jogging, Badminton, Hiking. %\href{https://ryannnice.github.io//}{Renyuan's Homepage}

\end{document}
